In einem Labor liegt auf einem isolierenden Tisch eine Kupferleitung mit einem Querschnitt von \(\SI{4,00}{\milli\meter\squared}\) und einer Länge von \(\SI{2,00}{\meter}\). Die Leitung ist über Zuleitungen mit einer Spannungsquelle verbunden, so dass über der Kupferleitung eine Spannung von \(\SI{0,590}{\volt}\) abfällt; die Zuleitungen werden senkrecht nach unten durch den Tisch geführt, siehe Abbildung~\ref{leiterdeckefig}.

Parallel zu dieser Leitung hängt im Labor von der Decke an zwei identischen, vertikal angeordneten Federn eine weitere identische Kupferleitung. Auch diese hängende Leitung ist über Zuleitungen mit einer Spannungsquelle verbunden, so dass durch beide Leitungen Ströme gleichen Betrags in die gleiche Richtung fließen, siehe Abbildung~\ref{leiterdeckefig}. Die Zuleitungen zur hängenden Leitung werden senkrecht nach oben durch die Decke geführt; ihre Einflüsse auf den Aufbau werden vernachlässigt.

Ohne Stromfluss durch die Leitungen hängt die zweite Kupferleitung im statischen Gleichgewicht an den beiden Federn. Wird nun der Strom in beiden Leitungen eingeschaltet, so so wird jede Feder gegenüber dem stromlosen Gleichgewichtszustand um \(\SI{2,00}{\milli\meter}\) gedehnt. Im so gedehnten Zustand der Federn beträgt der vertikale Abstand der Leiter \(\SI{50,0}{\milli\meter}\). 

Berechnen Sie die Federkonstante, die jede der beiden Federn haben muss.